\section{A sufficient estimate for ordered logarithmic scales}
\label{sec:orderedscales}

Choi, Kadhe, and Puhm identify the persistence of the logarithmic
soft-factor tower beyond one loop as an open quantum test
\cite[Sec.~6]{longrange}. For massless scalar QED, the first test is
the coefficient of $\omega\log^2(\omega/\Lambda)$ in the
symmetry-governed single-soft-photon insertion, with hard scale
$\Lambda$. In their positive-helicity convention, define the
tree prefactor $S_{-1,i}=\sqrt{2}eQ_i/(z-z_i)$ and write
\[
 B_i=-\gamma_{\mathrm{em}}\eta_i
       \sum_{j\ne i}\frac{(\bar z-\bar z_j)Q_iQ_j}
                          {\bar z_i-\bar z_j}\,
         e^{-\partial_{\Delta_i}} .
\]
Here $z_i$ are distinct hard punctures, $Q_i$ and $\eta_i$ are
the charges and in/out signs in their convention,
$\gamma_{\mathrm{em}}=e^2/(4\pi^2)$, and the exponential
shifts the hard weight $\Delta_i$. Equation~(4.24) of
\cite{longrange} uses the normalized Mellin coefficient
in their Eq.~(4.11). In this normalization the energy
coefficient is preserved, since
\[
 \int_0^\Lambda \omega^{\Delta-1}\omega\log^2(\omega/\Lambda)\,d\omega
 =\frac{2\Lambda^{\Delta+1}}{(\Delta+1)^3}
 \qquad(\Re\Delta>-1)
\]
and multiplication by $(\Delta+1)^3/2$ extracts its leading
coefficient at $\Delta=-1$. The predicted contribution is
$S_{-1,i}B_i^2H/2$ on a fixed hard correlator $H$. To test
the contact coefficient at
puncture $i$, pair $\bar\partial^3$ of this term with a smooth
test supported near $z_i$ and away from the other hard punctures,
as motivated by the descendant in Eq.~(4.25). This projects the
product after composition. The following estimate gives
conditions on an independently calculated kernel sufficient
to obtain the $1/2$ in the projected leading coefficient.

\begin{proposition}[Ordered-scale leading coefficient]
\label{prop:orderedscales}
Let $X$ be a Banach space and $m\ge2$. For $T>0$ use the ordered domain
\[
 S_m(T)=\{(t_1,\ldots,t_m):0<t_1<\cdots<t_m<T\}.
\]
Suppose
\[
 F_m(T)=\int_{S_m(T)}K_T(t)\,dt+E_T,\qquad
 \|E_T\|\le C T^{m-1},
\]
where $K_T$ is Bochner integrable. Define $d_0=t_1$,
$d_j=t_{j+1}-t_j$ for $1\le j<m$, and $d_m=T-t_m$.
Suppose there are $A_m\in X$ and nonnegative
$w_0,\ldots,w_m\in L^1([0,\infty))$, independent of $T$, such that
\[
 \|K_T(t)-A_m\|\le\sum_{j=0}^m w_j(d_j)
 \quad\text{for almost every }t\in S_m(T).
\]
Then
\begin{equation}
 \left\|F_m(T)-\frac{T^m}{m!}A_m\right\|
 \le T^{m-1}\left(C+\frac1{(m-1)!}\sum_{j=0}^m\|w_j\|_1\right).
 \label{eq:orderedscale-error}
\end{equation}
\end{proposition}
\begin{proof}
The change from the ordered $t_j$ to the independent gaps
$d_0,\ldots,d_{m-1}$ has unit Jacobian and gives the simplex
of nonnegative gaps with sum at most $T$. The remaining gap
is $d_m=T-\sum_{j=0}^{m-1}d_j$. Any gap can be chosen as the
dependent coordinate, since the corresponding affine
substitution has Jacobian of absolute value one. At fixed
$d_j=u$, the other independent coordinates form a simplex
of dimension $m-1$ with volume $(T-u)^{m-1}/(m-1)!$.
Consequently
\begin{equation}
 \int_{S_m(T)}w_j(d_j)\,dt=
 \frac1{(m-1)!}\int_0^T (T-u)^{m-1}w_j(u)\,du
 \le\frac{T^{m-1}}{(m-1)!}\|w_j\|_1 .
 \label{eq:orderedscale-gap}
\end{equation}
Since $S_m(T)$ has volume $T^m/m!$, the constant part
integrates to $T^m A_m/m!$. Apply the integral norm
inequality to the difference and include the bound on $E_T$.
\end{proof}

The error power in \eqref{eq:orderedscale-error} is sharp.
With $E_T=0$ and $K_T-A_m=e^{-d_j}v$ for a unit vector $v$,
the integrated norm is asymptotic to $T^{m-1}/(m-1)!$ by
\eqref{eq:orderedscale-gap} and dominated convergence.
A slowly decaying choice $(1+d_j)^{-1}v$ instead produces
\[
 \frac{T^{m-1}}{(m-1)!}\big(\log(1+T)+O(1)\big).
\]
Indeed, the difference between the scaled integral in
\eqref{eq:orderedscale-gap} and $\int_0^T (1+u)^{-1}\,du$
is bounded using $0\le1-(1-u/T)^{m-1}\le (m-1)u/T$.
Thus pointwise decay at large gap does not guarantee the
error rate in \eqref{eq:orderedscale-error}. For $m=2$, a
constant additional kernel $v$ would give a contribution
$T^2v/2$, even with all one-loop data held fixed.

For the QED test, take $T=\log(\Lambda/\omega)$ and write
virtual scales as $\kappa_j=\omega e^{t_j}$, so
$d\kappa_j/\kappa_j=dt_j$. To apply
Proposition~\ref{prop:orderedscales}, one must extract
$F_2$ after dividing by the single soft-energy power
$\omega$ and projecting in a fixed regulator and
soft-limit prescription. Combining all contributing
orderings and normalizations, controlling collinear and
overlap subtractions, and bounding omitted regions by
$O(T)$ are part of deriving the ordered representation.
The well-separated kernel must then match the projected
$S_{-1,i}B_i^2H$, with its deviation controlled also at
comparable scales and the two endpoints. A derivation of
ordinary virtual infrared exponentiation, such as the
gravity calculation of Naculich and Schnitzer
\cite{naculichschnitzer}, concerns a virtual divergence
with different kinematics. The ordered representation
and kernel estimate for the radiative QED coefficient
remain to be established.
